%=====================================================================
\chapter{Reference Management and Bibliometrics}\label{ch:bibliometrics}
%=====================================================================
\locator{2}{Part~2 --- Finding and Synthesising Knowledge}

\begin{outcomes}
By the end of this chapter you will be able to:
\begin{tight}
  \item \textbf{Operate} an open-source reference manager as the single source
    of truth for a review.
  \item \textbf{Write} BibTeX entries that will still resolve in five years.
  \item \textbf{Distinguish} co-citation, bibliographic coupling and co-word
    analysis, and say what each reveals.
  \item \textbf{Produce} a field map from a Scopus or Web of Science export
    using open-source tools.
  \item \textbf{Judge} when a bibliometric analysis is a contribution and when
    it is a substitute for reading.
\end{tight}
\end{outcomes}

\begin{prereq}
\chref{slr} for the review workflow this tooling supports, and \chref{publication}
for the critique of citation metrics as evaluation --- this chapter uses them
as \emph{data} instead.
\end{prereq}

\begin{keyterms}
reference manager $\bullet$ BibTeX and biblatex $\bullet$ citation key
$\bullet$ DOI $\bullet$ bibliometrics $\bullet$ co-citation $\bullet$
bibliographic coupling $\bullet$ co-word analysis $\bullet$ co-authorship
network $\bullet$ science mapping
\end{keyterms}

\section{Never Type a Citation Again}

Hand-typed references are wrong. Not occasionally --- systematically. Volume
numbers drift, initials get dropped, a conference is recorded as a journal, and
the error propagates because the next author copies your reference list rather
than the source. In a review with 80 primary studies, hand management
guarantees a corrupted bibliography and hours lost to reformatting when you
change venue.

\begin{center}
\small
\begin{tabular}{@{}L{2.4cm}L{3.0cm}L{9.0cm}@{}}
\toprule
\textbf{Tool} & \textbf{Licence} & \textbf{Why you might choose it} \\
\midrule
Zotero   & Open source (AGPL)
         & Browser connector captures metadata and the PDF in one click; group
           libraries make two-screener review practical; exports BibTeX.
           The default recommendation \\
JabRef   & Open source (MIT)
         & Native BibTeX: the \texttt{.bib} file \emph{is} the database, so
           there is no export step to go stale. Best if you write in \LaTeX{}
           and want one file under version control \\
Mendeley & Proprietary
         & Widely used; owned by a publisher, with the data-ownership questions
           that raises \\
\bottomrule
\end{tabular}
\end{center}

\begin{conceptbox}[A workflow that survives a three-year degree]
\begin{tightnum}
  \item \textbf{One library, from day one.} Every paper you read goes in
    immediately, including ones you reject --- with a note saying why. In month
    20 you will not remember, and you will re-read it.
  \item \textbf{Collections per review stage.} \texttt{01-identified},
    \texttt{02-screened}, \texttt{03-included}, and \texttt{99-excluded} with
    the reason in each item's note. Your PRISMA counts then come out of the
    library rather than out of a spreadsheet you forgot to update.
  \item \textbf{Fix metadata on import.} Automatic capture is good, not
    perfect. Check the DOI, the venue and the year while the paper is in front
    of you.
  \item \textbf{Export \texttt{refs.bib} into the repository} and commit it
    with the thesis. The \texttt{refs.bib} shipped with this handout is exactly
    that file for this course.
  \item \textbf{Back it up somewhere that is not your laptop.}
\end{tightnum}
\end{conceptbox}

\subsection{BibTeX entries that keep working}

\begin{center}
\small
\begin{tabular}{@{}L{3.0cm}L{11.3cm}@{}}
\toprule
\textbf{Rule} & \textbf{Why} \\
\midrule
Stable keys      & \texttt{author\-year\-firstword}, lowercase. Never renumber;
                   a key is an identifier, not a position \\
DOI on everything & The only durable link. URLs rot; DOIs resolve \\
Correct entry type & \texttt{@inproceedings} for conference papers,
                   \texttt{@article} for journals. Getting this wrong misformats
                   the whole reference \\
Protect capitals & \verb|{PRISMA}|, \verb|{IEEE}|, \verb|{FAIR}| --- otherwise
                   many styles lowercase them \\
Full author lists & Let the style truncate to \emph{et al.}, not the data.
                   Truncating in the \texttt{.bib} loses information permanently \\
Archived URL + \texttt{urldate} & For anything grey (\chref{greylit}) \\
\bottomrule
\end{tabular}
\end{center}

\begin{pitfall}[The two errors that reach print]
\textbf{Citing something you have not read.} You copy a reference from another
paper's list; that paper misdescribed it; now you have too. If you cite it, open
it. If you genuinely cannot obtain it, cite it as reported in the source you
read --- ``[8], as cited in [23]'' --- which is honest and rare enough to be
noticed favourably.

\medskip
\textbf{Letting the manager write the sentence.} A reference manager gets the
formatting right and the \emph{meaning} not at all. It cannot tell you that the
paper you are citing for a claim actually concluded the opposite under
conditions you did not read. \chref{appraisal} is about that.
\end{pitfall}

\section{Bibliometrics as a Research Method}

\chref{publication} argued that citation metrics are poor instruments for
\emph{evaluating} people. That is a separate question from using citation data
as \emph{evidence about a field} --- which is legitimate, well-established, and
occasionally a contribution in its own right.

\begin{center}
\small
\begin{longtable}{@{}L{3.2cm}L{5.0cm}L{6.0cm}@{}}
\toprule
\textbf{Technique} & \textbf{Two items are linked when} & \textbf{What it shows} \\
\midrule
\endfirsthead
\toprule
\textbf{Technique} & \textbf{Two items are linked when} & \textbf{What it shows} \\
\midrule
\endhead
\bottomrule
\endfoot
Co-citation
  & Both are cited by the same later paper
  & The intellectual base a community shares. Backward-looking; identifies
    foundational works and schools of thought \\
Bibliographic coupling
  & Both cite the same earlier paper
  & Current research fronts. Forward-looking; groups work being done now,
    before citation has accumulated \\
Co-word / co-occurrence
  & Two terms appear together in titles, abstracts or keywords
  & The conceptual structure and its vocabulary; useful for finding search
    terms you missed (\chref{slr}) \\
Co-authorship
  & Two authors share a paper
  & Collaboration structure; which groups exist and whether they talk \\
\end{longtable}
\end{center}

\begin{conceptbox}[The distinction that matters most]
Co-citation and bibliographic coupling look similar and answer opposite
questions.

\medskip
Co-citation needs \emph{time}: two papers only become co-cited once enough
later work has cited both. It therefore describes the settled past, and a
brand-new paper cannot appear in a co-citation map at all.

\medskip
Bibliographic coupling needs \emph{no} time: two papers published last month
that cite the same twenty references are strongly coupled immediately. It is
the right tool for ``what is happening now'', which is usually what a scholar
choosing a topic actually wants to know.
\end{conceptbox}

\subsection{Doing it with open-source tools}

\begin{center}
\small
\begin{tabular}{@{}L{2.6cm}L{3.2cm}L{8.6cm}@{}}
\toprule
\textbf{Tool} & \textbf{What it is} & \textbf{Use} \\
\midrule
VOSviewer  & Free (not open source)
           & Point-and-click maps from Scopus, WoS, Dimensions, Lens exports.
             The fastest route to a publishable map \\
bibliometrix / biblioshiny & R package, open source
           & Scriptable and reproducible, which VOSviewer is not.
             \texttt{biblioshiny()} gives a browser interface over it \\
CiteSpace  & Free
           & Emphasises temporal evolution and burst detection \\
\bottomrule
\end{tabular}
\end{center}

\begin{lstlisting}[language=Rlang, caption={A minimal reproducible field map. The
whole point of scripting it is that the figure in your thesis can be
regenerated when you add 40 more records.}]
library(bibliometrix)

# Export from Scopus as BibTeX, or from Web of Science as plain text.
M <- convert2df("scopus_export.bib", dbsource = "scopus", format = "bibtex")

results <- biblioAnalysis(M)
summary(results, k = 20)          # top authors, sources, countries, growth

# Conceptual structure: which terms occur together?
CS <- conceptualStructure(M, field = "ID", method = "MCA",
                          minDegree = 8, clust = "auto")

# Research fronts: bibliographic coupling of documents.
NetMatrix <- biblioNetwork(M, analysis = "coupling",
                           network = "documents", sep = ";")
networkPlot(NetMatrix, n = 40, type = "fruchterman",
            labelsize = 0.7, cluster = "louvain")
\end{lstlisting}

\begin{alertbox}[A map is a hypothesis-generator, not a finding]
A co-word map will show clusters. Clusters are produced by any clustering
algorithm on any data, including noise, and their number depends on parameters
you chose. The map tells you where to look; it does not tell you what is there.

\medskip
Every cluster you report should be validated by \emph{reading} several papers
in it and confirming that the label you gave it is the label they would
recognise. A bibliometric section that reports six clusters without
demonstrating that anybody read them is the most sophisticated possible way of
not doing a literature review.
\end{alertbox}

\begin{conceptbox}[When bibliometrics is itself the contribution]
A bibliometric study can be a legitimate contribution --- of the ``survey /
review'' type in \chref{contribution} --- when it answers a question about the
field that could not be answered by reading. Has research on this topic shifted
from academia to industry? Do the two communities studying this problem cite
each other? Has the geographic distribution of authorship changed as the
problem became commercially important?

\medskip
Those are real questions with real answers. What is not a contribution is a
descriptive count of papers per year with a growth curve, which is true of
almost every topic and tells nobody anything.
\end{conceptbox}

\begin{traditions}{mapping a field}
\tradEMP{Bibliometrics is an observational study of the literature: it has a
population, a sampling frame with known coverage limits, and confounders ---
notably that database coverage differs systematically by region and venue.}
\tradDSR{Maps the artefact space: which classes of solution have been built,
by whom, and where practitioner and academic communities diverge.}
\tradQUAL{Treats the corpus as documents to be read closely; the map is a
sampling device for choosing which documents deserve that reading.}
\tradFORM{Citation graphs are objects in their own right, and the algorithms
that cluster them have properties worth stating --- resolution limits,
parameter sensitivity, instability across runs.}
\end{traditions}

\begin{lab}[Map your own field]
\begin{tightnum}
  \item Run your \chref{slr} search on Scopus or Web of Science and export the
    full records, including references, for at least 200 documents.
  \item Load them in \texttt{bibliometrix} and produce three artefacts: a
    co-word map, a bibliographic coupling map of documents, and a table of the
    ten most locally cited papers.
  \item Read the abstracts of three papers from each cluster and name the
    cluster yourself. Compare your labels with the automatic ones.
  \item Extract from the co-word map at least two terms that were \emph{not} in
    your search string. Add them and re-run the search from \chref{slr}.
\end{tightnum}

\medskip
Step 4 is the reason this chapter sits inside Part~II rather than being an
optional extra: the map repairs the search.
\end{lab}

\begin{checkpoint}
\begin{tightnum}
  \item Why can a paper published this month appear in a bibliographic coupling
    map but not in a co-citation map?
  \item Give two reasons a reference list should live in a manager rather than
    a document, beyond convenience.
  \item A thesis reports a co-word analysis identifying five research themes.
    What single piece of evidence would most increase your confidence in those
    five labels?
  \item Scopus and Web of Science have different coverage. How would that
    difference bias a bibliometric study of research output by country?
\end{tightnum}
\end{checkpoint}

\begin{chaptersummary}
\begin{tight}
  \item Use an open-source reference manager from the first paper, with
    collections mirroring the review stages, so PRISMA counts fall out of the
    library.
  \item Stable keys, DOIs, correct entry types, protected capitals, full author
    lists, archived URLs for grey sources.
  \item Cite only what you have opened. If you truly cannot, say ``as cited
    in''.
  \item Co-citation shows the settled intellectual base; bibliographic coupling
    shows current fronts and works immediately; co-word shows vocabulary and
    repairs your search string.
  \item \texttt{bibliometrix} is scriptable and therefore reproducible;
    VOSviewer is faster but its figures cannot be regenerated automatically.
  \item A map generates hypotheses. Validate every cluster label by reading, or
    the analysis is a way of avoiding the literature rather than reviewing it.
\end{tight}
\end{chaptersummary}

\begin{reviewq}
  \item \chref{publication} argued citation counts are poor measures of
    individual quality, yet this chapter uses citation data as evidence about
    fields. Is this inconsistent? Set out the conditions under which aggregate
    citation analysis is sound even though individual-level use is not.
  \item Database coverage is not neutral: Scopus and Web of Science index
    journals from some regions and languages far more completely than others.
    What does this imply for a bibliometric claim about research productivity
    in South Asia?
  \item Design a bibliometric study that would answer a question about your own
    field that could not be answered by reading fifty papers. State the
    question, the data, the technique, and the threat to validity you would
    have to address.
  \item Clustering a citation network requires parameter choices that change
    the number of clusters. What reporting practice would make a bibliometric
    result as trustworthy as a well-reported statistical one?
\end{reviewq}
